Dans l'espace muni d'un repère orthonormé \Oijk, on considère les points
\[\mathrm{A}(1 ; 2 ; 3),\quad \mathrm{B}(-1 ; 3 ; 1),\quad \mathrm{C}(2 ; 1 ; 6) \quad \text { et } \quad \mathrm{D}(3 ;-2 ;-1).\]

	\begin{enumerate}
		\item \begin{enumerate}
			\item Montrer que les points A, B et C définissent un plan.

			\item Montrer que le vecteur $\vect{n}\begin{pmatrix} 1 \\ 4 \\ 1 \end{pmatrix}$ est normal au plan (ABC).

			\item  En déduire une équation cartésienne du plan (ABC).
		\end{enumerate}

		\item \begin{enumerate}
			\item Déterminer une équation paramétrique de la droite $(d)$, perpendiculaire au plan (ABC) et passant par le point D.

			\item Déterminer les coordonnées du point H qui est le projeté orthogonal du point D sur le plan (ABC).

			\item En déduire que la distance du point D au plan (ABC) est égale à $3\sqrt{2}$.
		\end{enumerate}

		\item \begin{enumerate}
			\item Montrer que
			\[\cos\left(\,\widehat{\mathrm{BAC}}\,\right)=-\dfrac{3\sqrt{11}}{11}\]

			\item En déduire la valeur exacte de $\sin\left(\,\widehat{\mathrm{BAC}}\,\right)$.

			\item Montrer que l'aire du triangle ABC vaut $\dfrac{3\sqrt{2}}{2}$.
		\end{enumerate}

		\item Déterminer le volume du tétraèdre ABCD.

		On rappelle que le volume $\mathcal{V}$ d'un tétraèdre est donné par la formule suivante :
		\[\mathcal{V}=\dfrac{1}{3}\times\mathcal{B}\times h\, ,\quad \text{ où } \mathcal{B}\text{ est l'aire d'une base et } h\text{ la hauteur qui lui est associée.}\]
	\end{enumerate}
